\section{Lineare Abhängigkeit}
\label{sec:vector_algebra_linear_dependence}

In den vorangegangenen Abschnitten haben wir Vektoren addiert, skaliert und
als Richtungen und Verschiebungen interpretiert. Nun fragen wir, ob mehrere
Vektoren tatsächlich neue Richtungen liefern oder ob sich einige von ihnen aus
anderen zusammensetzen und damit überflüssig sind.

\bigskip

Um diese Fragestellung besser zu verstehen, betrachten wir drei Vektoren
$\vec a,\vec b,\vec c\in\R^3$. Spannen $\vec a$ und $\vec b$ eine Ebene auf,
so kann der Vektor $\vec c$ die beiden anderen Vektoren nur kompensieren,
wenn er ebenfalls in dieser Ebene liegt. Mathematisch beschreiben wir die
Kompensationsfrage durch

\noindent
\begin{minipage}[c]{0.55\linewidth}
  \[
    \alpha\vec a+\beta\vec b+\gamma\vec c=\vec 0.
  \]
  Diese Gleichung beschreibt die Kompensation des dargestellten Vektorzuges:
  Die drei skalierten Vektoren bilden zusammen einen geschlossenen Zug.
\end{minipage}\hfill
\begin{minipage}[c]{0.40\linewidth}
  \centering
  \CachedTikzFigure{content/resources/chapters/vector_algebra/tikz/linear_dependence_construction_and_compensation.tex}
\end{minipage}

Schreiben wir die Vektoren in Koordinaten aus, wird aus der
Kompensationsgleichung
\[
  \alpha\begin{pmatrix}a_1\\a_2\\a_3\end{pmatrix}
  +\beta\begin{pmatrix}b_1\\b_2\\b_3\end{pmatrix}
  +\gamma\begin{pmatrix}c_1\\c_2\\c_3\end{pmatrix}
  =\begin{pmatrix}0\\0\\0\end{pmatrix}.
\]
Die Vektorsumme ist genau dann null, wenn jede ihrer Komponenten null ist.
Damit erkennen wir: Es handelt sich um ein homogenes lineares Gleichungssystem.
In Koordinaten lautet es:
\[
  \begin{aligned}
    a_1\alpha+b_1\beta+c_1\gamma&=0,\\
    a_2\alpha+b_2\beta+c_2\gamma&=0,\\
    a_3\alpha+b_3\beta+c_3\gamma&=0.
  \end{aligned}
  \qquad\Longleftrightarrow\qquad
  \left[
    \begin{array}{ccc|c}
      a_1&b_1&c_1&0\\
      a_2&b_2&c_2&0\\
      a_3&b_3&c_3&0
    \end{array}
  \right].
\]
Da ein homogenes Gleichungssystem immer die triviale Lösung
$(\alpha,\beta,\gamma)=(0,0,0)$ besitzt und uns deren existenz nichts
über die beteiligten Vektoren aussagt, schließen wir diese aus und suchen
ausschließlich nichttriviale Lösungen:
\[
  (\alpha,\beta,\gamma)\ne(0,0,0).
\]

Damit gilt: Besitzt das Gleichungssystem eine ...
\begin{itemize}[leftmargin=1.6em,itemsep=0.25em]
  \item eine nichttriviale Lösung, ist eine
    gegenseitige Kompensation möglich: Die Vektoren sind linear abhängig.
  \item nur die triviale Lösung, ist keine gegenseitige Kompensation
    möglich: Die Vektoren sind linear unabhängig.
\end{itemize}

Jetzt übertragen wir dieses Konzept auf mehr Vektoren und höhere Dimensionen.

\begin{DefinitionBox}{Lineare Abhängigkeit und Unabhängigkeit}{vector_algebra_linear_dependence}
  Die Vektoren $\vec v_1,\dots,\vec v_k\in\R^n$ heißen \emph{linear abhängig}
  \EnglishTerm{linearly dependent}, wenn das homogene Gleichungssystem
  \[
    \lambda_1\vec v_1+\dots+\lambda_k\vec v_k=\vec 0
  \]
  eine nichttriviale Lösung besitzt, also mindestens ein
  $\lambda_i\ne0$ ist. Besitzt es nur die triviale Lösung
  $\lambda_1=\dots=\lambda_k=0$, heißen die Vektoren \emph{linear unabhängig}
  \EnglishTerm{linearly independent}.

  Da ein homogenes LGS nicht ohne Lösung sein kann, ergibt sich damit die
  folgende Zuordnung:
  \begin{center}
    \CachedTikzFigure{content/resources/chapters/vector_algebra/tikz/linear_dependence_solution_types.tex}
  \end{center}
\end{DefinitionBox}

\noindent
\begin{minipage}{\linewidth}
Die Frage nach linearer Abhängigkeit oder Unabhängigkeit hängt somit von
der Lösungsmenge des zugehörigen homogenen LGS ab. 

Stimmen Anzahl der Vektoren
und Dimension überein, ist die zugehörige Matrix quadratisch. Genau dann kann
die Determinante die Entscheidung abkürzen: Statt das homogene LGS vollständig
mit dem Gauß-Verfahren zu lösen, genügt die Berechnung einer Determinante.
\end{minipage}

\begin{KeyPointBox}{Determinantenkriterium für Vektoren}{vector_algebra_linear_dependence_determinant_criterion}
  Für
  \[
    \underbrace{\vec v_1,\dots,\vec v_n}_{\#\text{ Vektoren}=n}
    \qquad\text{mit}\qquad
    \vec v_i\in
    \underbrace{\R^n}_{\#\text{ Dimensionen}=n}
  \]
  gilt $\#\text{ Vektoren}=\#\text{ Dimensionen}$.
  Wir schreiben die Vektoren als Spalten und erhalten die quadratische
  Determinante
  \[
    \det(\vec v_1,\vec v_2,\dots,\vec v_n).
  \]
  Für das homogene LGS aus
  Definition~\ref{def:vector_algebra_linear_dependence} gilt
  \[
    \det(\vec v_1,\dots,\vec v_n)\ne0
    \;\Longleftrightarrow\;
    \begin{gathered}
      \lambda_1\vec v_1+\dots+\lambda_n\vec v_n=\vec 0\\
      \text{hat nur die triviale Lösung}
    \end{gathered}
    \;\Longleftrightarrow\;
    \begin{gathered}
      \vec v_1,\dots,\vec v_n\\
      \text{sind linear }\Cloze{3.7cm}{\text{unabhängig}}
    \end{gathered}.
  \]

  Durch Negation ergibt sich
  \[
    \det(\vec v_1,\dots,\vec v_n)=0
    \;\Longleftrightarrow\;
    \begin{gathered}
      \lambda_1\vec v_1+\dots+\lambda_n\vec v_n=\vec 0\\
      \text{hat unendlich viele Lösungen}
    \end{gathered}
    \;\Longleftrightarrow\;
    \begin{gathered}
      \vec v_1,\dots,\vec v_n\\
      \text{sind linear }\Cloze{3.7cm}{\text{abhängig}}
    \end{gathered}.
  \]
  Eine Determinante ist nur für quadratische Matrizen definiert. Das Kriterium
  ist daher genau dann anwendbar, wenn
  $\#\text{ Vektoren}=\#\text{ Dimensionen}$ gilt. In diesem Fall ersetzt die
  Determinantenberechnung die vollständige Lösung des homogenen LGS und spart
  häufig das Gauß-Verfahren. Stimmen beide Anzahlen nicht überein, untersuchen
  wir das homogene LGS mit dem Gauß-Verfahren.
\end{KeyPointBox}

\begin{KeyPointBox}{Mehr Vektoren als Dimensionen}{vector_algebra_linear_dependence_more_vectors_than_dimensions}
  Für $\vec v_1,\dots,\vec v_k\in\R^n$ mit $k>n$ hat das zugehörige
  homogene LGS $k$ Unbekannte, aber nur $n$ Gleichungen. Seine erweiterte
  Koeffizientenmatrix kann beispielsweise die Zeilenstufenform
  \[
    [A\mid\vec0]\sim
    \left[
      \begin{array}{cccc|c}
        \ast&\ast&\ast&\ast&0\\
        0&\ast&\ast&\ast&0\\
        0&0&\ast&\ast&0
      \end{array}
    \right]
    \qquad
    \begin{gathered}
      3~\text{Gleichungen, }4~\text{Unbekannte}\\[-0.1em]
      \text{mindestens eine freie Variable}
    \end{gathered}
  \]
  besitzen. Allgemein gilt
  \[
    \operatorname{rang}(A)=\operatorname{rang}([A\mid\vec0])
    \le n<k=\#\text{ Unbekannte}.
  \]
  Das homogene LGS besitzt daher unendlich viele, insbesondere nichttriviale
  Lösungen. Die Vektoren sind linear abhängig.
\end{KeyPointBox}

Anhand des folgenden Ablaufdiagramms kann man sich im jeweiligen Fall
orientieren, wie die lineare Abhängigkeit oder Unabhängigkeit nachzuweisen ist.
\begin{center}
  \CachedTikzFigure{content/resources/chapters/vector_algebra/tikz/linear_dependence_method_choice.tex}
\end{center}

\begin{KeyPointBox}{Geometrische Interpretation}{vector_algebra_linear_dependence_geometric_interpretation}
  \begin{itemize}[leftmargin=1.6em,itemsep=0.55em,topsep=0.2em]
    \item Für zwei Vektoren in $\R^2$ gilt:
      \[
        \vec a,\vec b\text{ sind linear abhängig}
        \quad\Longleftrightarrow\quad
        \vec a\parallel\vec b.
      \]
      Sie spannen dann keine Fläche auf.
    \item Für drei Vektoren in $\R^3$ gilt:
      \[
        S=(\vec a\times\vec b)\cdot\vec c
        =\det(\vec a,\vec b,\vec c),
        \qquad
        V=|S|.
      \]
      Die Vektoren sind genau dann linear unabhängig, wenn sie ein Volumen
      aufspannen, also $V>0$ gilt.
  \end{itemize}
\end{KeyPointBox}

\begin{ExampleBox}{Lineare Abhängigkeit und Unabhängigkeit prüfen}{vector_algebra_linear_independence_determinant}
  Wir betrachten drei typische Fälle.
  \begin{enumerate}[label=\alph*),leftmargin=*,itemsep=0.8em]
    \item \textbf{Drei Vektoren in $\R^3$:}
      \[
        \vec a=\begin{pmatrix}2\\1\\1\end{pmatrix},\qquad
        \vec b=\begin{pmatrix}2\\3\\4\end{pmatrix},\qquad
        \vec c=\begin{pmatrix}7\\9\\11\end{pmatrix}.
      \]
      Die Determinante mit den Vektoren als Spalten ergibt
      \[
        \det(\vec a,\vec b,\vec c)
        =\det\!\begin{pmatrix}
          2&2&7\\
          1&3&9\\
          1&4&11
        \end{pmatrix}
        =2(33-36)-2(11-9)+7(4-3)
        =\Cloze{1.4cm}{-3}.
      \]
      Die drei Vektoren sind daher linear \Cloze{3.8cm}{unabhängig}.

    \item \textbf{Drei Vektoren in $\R^2$:}
      Für
      \[
        \vec a=\begin{pmatrix}1\\0\end{pmatrix},\qquad
        \vec b=\begin{pmatrix}0\\1\end{pmatrix},\qquad
        \vec c=\begin{pmatrix}1\\1\end{pmatrix}
      \]
      gilt $k=3>2=n$. Nach Merke~\ref{key:vector_algebra_linear_dependence_more_vectors_than_dimensions}
      sind die Vektoren daher linear \Cloze{3.4cm}{abhängig}.

    \item \textbf{Zwei Vektoren in $\R^3$:}
      Für
      \[
        \vec a=\begin{pmatrix}1\\0\\0\end{pmatrix},\qquad
        \vec b=\begin{pmatrix}0\\1\\0\end{pmatrix}
      \]
      verwenden wir das Gauß-Verfahren für das homogene LGS:
      \[
        \alpha\vec a+\beta\vec b=\vec0
        \quad\Longleftrightarrow\quad
        [A\mid\vec0]
        =\left[
          \begin{array}{@{}rr|r@{}}
            1&0&0\\
            0&1&0\\
            0&0&0
          \end{array}
        \right].
      \]
      Die Matrix liegt bereits in Zeilenstufenform vor. Damit gilt
      \[
        \operatorname{rang}(A)
        =\operatorname{rang}([A\mid\vec0])
        =\Cloze{1cm}{2}
        =\#\text{ Unbekannte}.
      \]
      Das LGS besitzt nur die triviale Lösung
      $(\alpha,\beta)=\Cloze{2.5cm}{(0,0)}$; die Vektoren sind daher linear
      \Cloze{3.8cm}{unabhängig}.
  \end{enumerate}
\end{ExampleBox}
